\documentclass[a4paper,12pt]{article}
\usepackage[margin=1.5cm]{geometry}
\usepackage{amsmath, amssymb}
\usepackage{array}
\usepackage{multicol}

\usepackage[spanish, bidi=basic, provide=*]{babel}


\usepackage{amsmath}
\usepackage{amssymb}
\usepackage{booktabs}
\usepackage{xcolor}
\usepackage{tcolorbox}
\usepackage{enumitem}
\usepackage{hyperref}

\definecolor{primary}{RGB}{31, 78, 121}
\definecolor{secondary}{RGB}{220, 53, 69}
\definecolor{bglight}{RGB}{245, 247, 250}

\hypersetup{
    colorlinks=true,
    linkcolor=primary,
    urlcolor=primary
}

\setlist[enumerate]{label=\textbf{\arabic*.--}, leftmargin=*}

\begin{document}

\begin{center}
    \Large\textbf{\color{primary} TEMA 2: POTENCIAS Y RAÍCES}
\end{center}

\begin{enumerate}

    \item \textbf{Reduce a una sola potencia:}
        \begin{enumerate}[label=\alph*)]
            \item $2^{3} \cdot 2 \cdot \left(\dfrac{2^{3} \cdot 2}{2^{4} \cdot 2^{2}}\right)$
            \item $16^{6} : (8^{5} \cdot 4^{2})$
            \item $\dfrac{2^{5} \cdot 3^{2} \cdot 4^{-1}}{2^{3} \cdot 9^{-1}}$
            \item $\left(\dfrac{1}{2}\right)^{3} : \left(\dfrac{1}{4}\right)^{2}$
        \end{enumerate}

    \item \textbf{Simplifica empleando las propiedades de las potencias, y expresa el resultado con exponentes positivos:}
        \begin{enumerate}[label=\alph*)]
            \item $\dfrac{3^{-2} \cdot (-7)^{2} \cdot 3 \cdot 7^{-4} \cdot (-3)^{5}}{(-7)^{3} \cdot 3^{-1} \cdot 7^{-5} \cdot 3^{4}}$
            \item $\dfrac{2^{3} \cdot 8^{-3} \cdot 12^{-1} \cdot (-3)^{-2}}{6^{2} \cdot 16^{-2} \cdot 3^{-3}}$
        \end{enumerate}

    \item \textbf{Simplifica empleando las propiedades de las potencias:}
        \begin{enumerate}[label=\alph*)]
            \item $\dfrac{[(2^{3}:3^{3}):(3^{-2}:2^{-2})]:6^{-2}}{[(4^{3}:7^{3}):3^{3}]:(6^{-1}:6^{-4})}$
            \item $\dfrac{9^{2} \cdot 3^{-3} \cdot 25}{125 \cdot 81}$
            \item $\dfrac{9^{5} \cdot 3^{-3} \cdot 25^{2}}{125 \cdot 27^{3}}$
            \item $\dfrac{81^{3} \cdot 3^{-3} \cdot 25^{-2}}{125^{-3} \cdot 27^{3}}$
        \end{enumerate}

    \item \textbf{Los terrenos de dos parcelas miden $3^{8}$ y $3^{4}$ metros cuadrados, respectivamente. Mohamed duda si la primera parcela es el doble que la segunda o no. De no ser doble, ¿cuántas veces es mayor la primera que la segunda?}

    \item \textbf{La masa de la Luna es de $7,34 \cdot 10^{22}$ kg, y la de la Tierra, de $5,98 \cdot 10^{24}$ kg. ¿A cuántas Lunas equivale la masa de la Tierra?}

    \item \textbf{Si el área del huerto cuadrado de mi abuelo es la mitad que el de mi tío, que tiene $200\text{ m}^{2}$, ¿cuánto mide el lado del huerto de mi abuelo?}

    \item \textbf{La unidad de memoria de un ordenador es el byte. Un kilobyte (kB) son $2^{10} = 1.024$ bytes, un megabyte (MB) son $2^{10}\text{ kB}$, un gigabyte (GB) equivale a $2^{10}\text{ MB}$ y un terabyte (TB) equivale a $2^{10}\text{ GB}$. Expresa en forma de potencia y en notación científica los bytes que tiene mi disco duro extraíble de $36\text{ TB}$.}

    \item \textbf{Queremos construir un cubo de cartón cuyo volumen sea $6\text{ m}^{3}$. ¿Qué superficie de cartón se necesita? Expresa el resultado en forma radical.}

    \item \textbf{Extrae todos los factores posibles de los siguientes radicales:}
        \begin{enumerate}[label=\alph*)]
            \item $\sqrt{81 \cdot a^{5} \cdot b \cdot c^{6}}$
            \item $\sqrt[3]{125 \cdot a^{9} \cdot b^{17} \cdot c^{25}}$
        \end{enumerate}

    \item \textbf{Efectúa las siguientes multiplicaciones y divisiones con radicales:}
        \begin{enumerate}[label=\alph*)]
            \item $14\sqrt{12} : 7\sqrt{3}$
            \item $\sqrt{2} \cdot \sqrt{8}$
            \item $\sqrt[3]{9} \cdot \sqrt[3]{3}$
            \item $\sqrt[3]{2} \cdot \sqrt{3}$
            \item $\sqrt[3]{2} \cdot \sqrt[3]{32}$
        \end{enumerate}

    \item \textbf{Extrae los factores que se puedan de los siguientes radicales:}
        \begin{enumerate}[label=\alph*)]
            \item $\sqrt[3]{\dfrac{8}{729} b^{5} c^{7} m^{14}}$
            \item $\sqrt{\dfrac{81}{32} z^{6} y^{7} x^{17}}$
            \item $\sqrt[3]{\dfrac{216}{343} m^{12} b^{15} c}$
            \item $\sqrt[4]{\dfrac{1}{243} b^{7} m^{45} c^{18} x^{5}}$
        \end{enumerate}

    \item \textbf{Considera las potencias $2^{-2}$, $2^{-3}$ y $2^{-5}$:}
        \begin{enumerate}[label=\alph*)]
            \item ¿Cuál es la mayor?
            \item ¿Cómo varía el valor de la potencia a medida que aumenta el valor absoluto del exponente negativo?
            \item Responde a las cuestiones anteriores para $0,7^{-3}$, $0,7^{-4}$ y $0,7^{-5}$.
        \end{enumerate}

    \item \textbf{Efectúa las siguientes sumas y restas de radicales:}
        \begin{enumerate}[label=\alph*)]
            \item $2\sqrt{20} + 4\sqrt{80} - 5\sqrt{180} + 3\sqrt{125}$
            \item $\dfrac{1}{4}\sqrt{128} + 6\sqrt{512} - \dfrac{1}{2}\sqrt{32} - 3\sqrt{98}$
            \item $\dfrac{3}{2}\sqrt{45} - \dfrac{\sqrt{20}}{3} + 4\sqrt{125} - 2\sqrt{5}$
        \end{enumerate}

    \item \textbf{Calcula los valores de $a, b, c, d$ en las siguientes igualdades:}
        \begin{enumerate}[label=\alph*)]
            \item $\sqrt{10^{4} \cdot 14^{6} \cdot 81^{12}} = 2^{a} \cdot 3^{b} \cdot 5^{c} \cdot 7^{d}$
            \item $\sqrt[3]{100^{9} \cdot 98^{9} \cdot 81^{15}} = 2^{a} \cdot 3^{b} \cdot 5^{c} \cdot 7^{d}$
        \end{enumerate}

    \item Efectua las operaciones con notación científica:
        \begin{enumerate}
            \item $5,3 \cdot 10^11 - 1,2 \cdot 10^12 + 7,2 \cdot 10^10 = $
            \item $4,2 \cdot 10^{-6} - 8,2 \cdot  10^{-7} + 1,8 \cdot 10^{-5} = $
            \item $(2,25 \cdot 10^22) \cdot (4 \cdot 10^{-15}) : (3 \cdot 10^{-3})=$
            \item $(1,4 \cdot 10^{-7}) : (5 \cdot 10^{-5}) =$ 
        \end{enumerate}

\end{enumerate}

\vspace{0.5cm}

\begin{center}
    \Large\textbf{\color{primary} TEMA 3: CÁLCULOS CON PORCENTAJES}
\end{center}

\begin{enumerate}
    \item \textbf{¿Qué porcentaje es?}
        \begin{enumerate}[label=\alph*)]
            \item El 40\% del 40\%
            \item El 25\% del 20\%
            \item El 30\% del 120\%
            \item El 150\% del 20\%
        \end{enumerate}

    \item \textbf{En la clase de 3º, han aprobado todas las asignaturas 24 estudiantes, que son el 75\% del total. ¿Cuántos estudiantes hay en la clase?}

    \item \textbf{En un partido de balonmano, una jugadora A ha marcado 2/5 de 30 intentos, otra jugadora B, 6 de 16 y una tercera jugadora C, el 36\% de 25 intentos.}
        \begin{enumerate}[label=\alph*)]
            \item ¿Cuántos goles ha marcado cada una?
            \item ¿Qué porcentaje de goles respecto al total de sus intentos ha marcado cada una?
        \end{enumerate}

    \item \textbf{El precio de un videojuego subió un 28\% y después lo rebajaron un 30\%. Si el precio inicial era de 58€ calcula el precio final.}

    \item \textbf{El precio de un ordenador portátil a principios de mes era de $640\text{ €}$. A lo largo del mes su precio ha subido un $15\%$. ¿Cuánto cuesta el ordenador ahora?}

    \item \textbf{En una ciudad había $45\,200$ estudiantes matriculados el curso pasado. Este año el número de matriculados ha disminuido un $8\%$. ¿Cuántos estudiantes hay matriculados ahora?}

    \item \textbf{El precio de una batidora, después de cargarle el $21\%$ de IVA, es de $72{,}6\text{ €}$. ¿Cuál es el precio antes de esos impuestos?}

    \item \textbf{Al estirar una goma elástica, su longitud aumenta un $30\%$ y en esta posición mide $104\text{ cm}$. ¿Cuánto mide sin estirar?}

    \item \textbf{En unas rebajas en las que se hace el $30\%$ de descuento, Roberto ha comprado una cámara por $50{,}4\text{ €}$. ¿Cuál era el precio inicial?}

    \item \textbf{Una cartera ha repartido el $36\%$ de las cartas que tenía. Aún le quedan $1\,184$. ¿Cuántas tenía antes de empezar el reparto?}

    \item \textbf{Un comercio aumenta el precio de sus productos un $30\%$ y después, pretendiendo dejarlo al precio inicial, los rebaja un $30\%$. ¿Es correcto?}
        \begin{enumerate}[label=\alph*)]
            \item Un ordenador costaba inicialmente $1\,000\text{ €}$. ¿Cuánto costará en cada paso del proceso?
            \item ¿Cuál es la variación porcentual que sufren los artículos respecto al precio inicial?
        \end{enumerate}

    \item \textbf{Un capital de $42\,000\text{ €}$ se deposita en un banco al $5\%$ de interés compuesto anual. ¿En cuánto se habrá convertido al cabo de un año, de dos y de tres años?}

\end{enumerate}

\end{document}